% common commands

\newcommand{\ifbooltrue}[2]{\ifthenelse{\boolean{#1}}{#2}{}}
\newcommand{\ifboolfalse}[2]{\ifthenelse{\boolean{#1}}{}{#2}}

\ifthenelse{\equal{\TwoSide}{true}}
{
    % TwoSide settings
    % Use default `\cleardoublepage'
}
{
    % OneSide settings
    \renewcommand{\cleardoublepage}{\clearpage}
}

\newcommand{\bful}[1]{{\bfseries\uline{#1}}}

% Commands to input body pages
\newcommand{\checkandinput}[2]
{
    % Directory structure:
    % graduate                        [Degree     ] [required        ]
    % |-- master/doctor               [GradLevel  ] [optional        ]
    %     |-- general/cs/math/...     [MajorFormat] [optional        ]
    %         |-- cover.tex/toc.tex/... [TeX Files  ] [template defined]
    %
    % Search from bottom-level to top-level:
    % E.g:
    %     With such dir:
    %     graduate
    %     |-- cover.tex
    %     |-- toc.tex
    %     |-- abstract.tex
    %     |-- master
    %         |-- toc.tex
    %         |-- math
    %             |-- abstract.tex
    %     The input files are:
    %         graduate/cover.tex
    %         graduate/master/toc.tex
    %         graduate/master/math/abstract.tex
    \IfFileExists{./#1/graduate/\GradLevel/\MajorFormat/#2}
    {
        \input{./#1/graduate/\GradLevel/\MajorFormat/#2}
    }
    {
        \IfFileExists{./#1/graduate/\GradLevel/#2}
        {
            \input{./#1/graduate/\GradLevel/#2}
        }
        {
            \input{./#1/graduate/#2}
        }
    }
}

\newcommand{\inputpage}[1]
{
    \checkandinput{page}{#1}
}

\newcommand{\inputbody}[1]
{
    \checkandinput{body}{#1}
}

\newcommand{\chapternonum}[1]
{
    \phantomsection
    \addcontentsline{toc}{chapter}{#1}
    \markboth{#1}{#1}
    % 无编号章标题（作者简历/摘要/致谢等）居中显示；
    % 用组限界，不影响编号正文各章（仍按 heading 的 \raggedright 格式）
    {
        \ctexset{chapter={format=\centering\zihao{-3}\bfseries}}
        \chapter*{#1}
    }
}

\newcommand{\sectionnonum}[1]
{
    \phantomsection
    \addcontentsline{toc}{section}{#1}
    \section*{#1}
}

% From: https://tex.stackexchange.com/questions/395856/switching-tocdepth-in-the-middle-of-a-document
\newcommand{\changelocaltocdepth}[1]{%
  \addtocontents{toc}{\protect\setcounter{tocdepth}{#1}}%
  \setcounter{tocdepth}{#1}%
}

%% ===== 修订标记：\rev{} / \revEN{}（统一开关 \ifshowrev）=====
% 显示/隐藏开关 \showrevtrue / \showrevfalse 在 zjuthesis.tex 中设置。
% 打开：显示标记；关闭：输出原文。
% 注意：\rev / \revEN 内不要放 \cite{}（会报错），引用放在外面：
%   \rev{文字部分}~\cite{key}。   % 正确
%   \rev{文字部分~\cite{key}}。   % 错误，勿用
\newif\ifshowrev

% 中文 / 混排版（内容可含 \ref{} 等宏）
\newcommand{\rev}[1]{%
  \ifshowrev
    \CJKunderanyline*{0.5ex}%
      {\color{yellow}\rule{0.5pt}{2.5ex}}%
      {#1}%
  \else
    #1%
  \fi
}

% 纯英文版（soul 背景高亮：跨行、保留词间空格；内容勿含 \cite/\ref 等宏）
\newcommand{\revEN}[1]{%
  \ifshowrev
    \hl{#1}%
  \else
    #1%
  \fi
}